\section{ISR y caché: hacer materialize del trabajo del backend en el edge}

En el comercio electrónico, el cuello de botella típico de Black Friday no es la red edge, sino que el origin/backend no soporta que cada usuario provoque un renderizado dinámico. Incremental Static Regeneration (ISR) hace \term{materialize} (materializa) el resultado de un renderizado como un artifact que se puede almacenar en caché, y lo vuelve a generar cuando vence la freshness policy o cuando ocurre un evento. Esto cambia el modelo de carga: el backend atiende los «cambios de contenido» y el edge atiende el «número de accesos».

\subsection{De la request amplification al backend shielding}

Si una página recibe $R$ solicitudes dentro de una freshness window, el renderizado dinámico tradicional requiere cerca de $R$ unidades de origin work; un ISR ideal solo necesita una generación y una revalidation, de modo que el trabajo del backend pasa de $O(R)$ a casi $O(1)$. Un sistema real también tiene que manejar stampede, recuperación ante fallas, cache regional, tag invalidation y stale content.

\lecturefigure{05-isr-cache-loop.png}{ISR usa materialization y revalidation para trasladar las lecturas repetidas del backend al edge.}{Subtítulos locales 05:55--07:20; documentación oficial de ISR de Next.js.}

\begin{knowledgebox}{Lectura de la figura: comparar la frecuencia de acceso con la frecuencia de cambio}
ISR es más adecuado para contenido que «se lee mucho, se escribe poco y admite una freshness explícita». Si los datos son distintos en cada solicitud, la tasa de aciertos de la cache es baja; si el contenido nunca puede estar desactualizado, se necesita un origin síncrono o una invalidation más estricta. La pregunta central no es «si se puede almacenar en caché», sino qué consistency contract permite el negocio.
\end{knowledgebox}

\begin{warningbox}{La caché no permite que el origin escale sin límite}
Los cache miss, las revalidation simultáneas, las solicitudes personalizadas y una gran cantidad de key distintas pueden seguir saturando el backend. La plataforma necesita request coalescing, stale-while-revalidate, protección de carga y un cache status visible; los desarrolladores también deben entender la semántica de la cache key y de la invalidation.
\end{warningbox}

\teachervoice{La forma de expresarse de Rauch tiene mucho sentido de ingeniería: que la nube sea «elastic» no significa que cualquier componente escale a la perfección. Uno de los valores de la plataforma es combinar automáticamente varios niveles de caché y la materialization, y convertir en capacidades del framework los patrones de Memcached, conjuntos de workers y escalamiento horizontal que antes los desarrolladores administraban a mano.}

\subsection{Resumen de la sección}

ISR desacopla la escala de accesos del cómputo del backend, de modo que una mayor parte del tráfico de la aplicación la absorben los artifact del edge. No consiste en volver estático todo, sino en administrar de forma explícita la freshness, la materialization y la revalidation.
